\chapter{实验研究类论文*}
\label{cha:thirdsection}


\section{引言（引言标题可选）}
\label{sec:method}
由于绪论中已对全文相关的研究背景和进展做了综述，因此，每章的引言中，请用1页左右的版面写一个导引，简要说明本章研究的背景或动机，以起到承上启下的作用，不宜太长。

引言的最后一段，说明本章的主要内容，如拟基于什么理论或方法，针对什么问题开展研究。注意：这里不能给结论。

若少于1个页面时，建议省略标题“引言”，直接在章的标题下写上几段话即可。

若学位论文属实验研究类论文，且论文所用的实验材料、仪器设备、实验方法基本一样，但应用不同，此时可考虑在第二章用整章的篇幅来描述，这样后面的章节就无需描重复描述相同内容，以避免冗余。

若实验性研究论文所用的材料与方法并非完全相同时，建议在各章中分别介绍材料与方法、实验结果、分析与讨论，最后是小结。




\section{分析与讨论}
论文排版时，无论因为图表等原因还是其它原因，除章与章的之间存在分页，空白处的地方不要太多。

讨论部分的关键在于：分析本章所得到的一些结果之间的关联性；所得到的结果与文献结果是否一致？分析产生一致或不一致结果的原因，以此体现论文的创新点。也可以指出本章中的不足，并分析原因及未来改进措施。

\section{本章小结}
本章主要介绍实验研究类的论文正文章节的框架结构。在每章的最后，都需要对该章的内容进行小结，不宜太长，建议1/2-2/3页版面较好。主要小结一下本章用什么理论或方法、做了什么事、得到的重要结果或结论。
